\section{The Change in Brief}
On October 9, 2026, WEPPpy changed standard low-severity forest initial random
roughness, \texttt{RRINIT}, from \textbf{4 to 6 cm} (0.04 to 0.06 m).
This corrects an inconsistency in the burned-forest parameter sequence.
Targeted tests support more sensible event-level sediment ordering, with
negligible runoff changes in the tested cases.\cite{adr}

This is a \textbf{WEPPpy parameterization change}, not a change to the
University of Idaho's \texttt{wepp-forest} equations or binaries. One shared
management-template field and four corresponding texture rows in the packaged
extended land/soil lookup changed. Ridge height (\texttt{rhinit}) remains
4 cm. Cover, soil, routing and other roughness defaults were not changed.

\begin{center}
\begin{tabular}{lrrrr}
\toprule
Standard class & Unburned & Low & Moderate & High\\
\midrule
Forest & 10 & \textbf{4 $\to$ 6} & 6 & 6\\
Young forest & 8 & \textbf{4 $\to$ 6} & 6 & 6\\
Shrub & 6 & 6 & 6 & 6\\
Tall grass & 2 & 2 & 2 & 2\\
\bottomrule
\end{tabular}
\end{center}
All values are centimeters of initial random roughness. Deciduous and mixed
forest share the forest values. Burned young forest uses the shared forest
templates. These are standard Disturbed mappings, \emph{not universal values
for every regional or custom mapping}. Shrub/grass burn remapping depends on
project options. Prescribed fire is separate and unchanged.

\textbf{Who is affected?} Projects preparing inputs from the shared low-severity
forest template, including the four forest labels above. Other mappings
selecting that template also inherit the revision. Different templates, custom
overrides and previously prepared inputs may not.

\clearpage
\section{Why Roughness Matters}
Initial random roughness is a modeled surface condition, not ridge height,
Manning's resistance coefficient, canopy height or litter depth. Measured litter
thickness should not simply be entered as RRINIT. The initial value also need
not persist: the model evolves an effective roughness, denoted here by $r$.

The retained code assessment traces roughness into depression storage,
hydraulic resistance and erosion-related calculations. Thus runoff, peak rates
and sediment can respond to this single input. The response depends on soil
state, slope and other inputs, not roughness alone.\cite{assessment}

One existing interrill-delivery factor explains the targeted sediment response:
\[
 f_r=\max(0,\min(1,1.14-23r)),\qquad r\text{ in meters}.
\]
Its cutoff is approximately 0.049565 m, or 4.96 cm. With effective roughness
at 4 cm the factor is 0.22; at 6 cm it is zero. Through the existing
particle-specific delivery calculation this suppresses the interrill
contribution. \textbf{It does not turn off all erosion.} Rill erosion, transport
and deposition remain separate processes. This describes an existing algorithm,
not a new physical law or a code change.

In the tested sandy-loam soil, effective roughness remained at the assigned
4 or 6 cm for all 36,525 days. In the other soils it declined below the cutoff,
eventually reaching the model's 0.6 cm floor. This explains why an input revision
can persist in one fixture but have no reported runoff or sediment effect in
others. It is not a general classification of field soils by texture.\cite{targeted}

\subsection{Why Not Lower All Values?}
The preceding 896-run sensitivity study tested four textures, six vegetation
labels, four severity states and two slope profiles. It compared class-specific
controls with values from 0.6 to 4 cm. Responses strongly depended on modeled
soil state. Reducing high-severity forest roughness from 6 to 0.6 cm on the
sandy-loam 20\% slope increased surface runoff volume by 17.89\% and sediment
delivery to about 2.91 times the control.\cite{sensitivity}

Those findings did not justify a blanket reduction. The adopted correction
uses the already established 6 cm burned-forest value for low severity, without
changing other parameters to compensate. Improved tested ordering supports the
consistency correction; it does not prove calibration against observations.
Lower sediment is not a degradation accepted for a neater parameter table.

\clearpage
\section{What the Targeted Comparison Shows}
The follow-up used 96 fresh canonical controls and 16 low-burn trials:
\textbf{112 successful simulations}. Four forest labels share burned templates;
identical results for a given soil are not independent replications.
Only RRINIT changed within each paired comparison.\cite{targeted}

All cases used the same frozen \texttt{wepp\_261009\_hill} binary and host runtime,
synthetic McKenzie Bridge climate (2000--2099; seed 26109), and an 87.9 m by
102.4 m profile with approximately 38.56\% mean grade. The follow-up did not
include the gentler profile. Surface runoff here is not routed streamflow;
hillslope sediment delivery is not channel erosion or watershed outlet sediment.

\begin{center}
\small
\begin{tabular}{lrrr}
\toprule
Soil fixture & Runoff depth & Sediment delivery & Maximum peak\\
 & mm / 100 yr & kg / 100 yr & m$^3$/s\\
\midrule
Clay loam & 41,620.0 $\to$ same & 111,737.85 $\to$ same & 0.11104 $\to$ same\\
Loam & 36,859.0 $\to$ same & 21,630.31 $\to$ same & 0.13427 $\to$ same\\
Sandy loam & 11,975.9 $\to$ same & 819.45 $\to$ 676.12 & 0.17035 $\to$ 0.17038\\
Silt loam & 25,505.6 $\to$ same & 8,824.65 $\to$ same & 0.12931 $\to$ same\\
\bottomrule
\end{tabular}
\end{center}
Arrows compare 4 with 6 cm. Runoff depth uses event-by-event (EBE) output;
sediment and peaks use hillslope PASS output. Unchanged values refer to the
retained reporting precision, not identical internal states.

\subsection{Sandy-Loam Water and Sediment}
Higher-precision PASS surface volume changed from 107,491.6793 to
107,491.6443 m$^3$: only $-0.035$ m$^3$ over 100 years ($-0.0000326$\%).
All 743 PASS event dates and 808 EBE dates were retained, with no new or lost
runoff-event dates. Peaks changed on 28 dates; the largest relative change was
approximately $-0.508$\%, and the 95th-percentile peak was unchanged.

PASS sediment decreased \textbf{17.49\%}, from 9.10 to 7.51 kg/ha/year over
the 0.900096 ha slope. Delivery decreased on 30 dates, increased on none, and
became zero on 24 previously positive-delivery dates. Annual totals changed in
25 of 100 years. This is persistent, event-selective behavior, not uniform
scaling or proof of physically zero sediment on those days.

The largest event mass decrease was 298.156 to 283.585 kg (4.89\%); another
event fell from 9.041 kg to zero. Rounded EBE delivery decreased 15.19\%, not
17.49\%; its kg/m measure differs from concentration-derived PASS mass.
The existing PASS surface-volume minus EBE depth-equivalent discrepancy changes
from $-302.9176$ to $-302.9526$ m$^3$. No accounting-closure repair is claimed.

\clearpage
\section{Rankings, Limits and User Actions}
\subsection{Improved Events, Not Universal Monotonicity}
Across the union of 3,465 sandy-loam mature-forest PASS event dates, there were
17 dates where low-severity sediment exceeded moderate-severity sediment at
4 cm, and none at 6 cm. Absent EVENT sediment is counted as zero in this
comparison. After revision, unburned $\leq$ low $\leq$ moderate $\leq$ high
delivery holds on every compared date.\cite{targeted}

Full-record rankings are distinct: runoff remained ordered in 23 of 24
soil/vegetation contexts, sediment in 24 of 24, and record maximum peak in
18 of 24. No full-record ordering failures were introduced or removed.
Remaining peak/runoff exceptions are not grounds to tune for a monotonic table.

Evidence covers one synthetic climate, four soil fixtures and a steep profile.
It does not establish field sediment accuracy, watershed outlet response,
subdaily timing parity, recovery through time or other-climate behavior.
The parameter study's 261009 identity is retained although that binary was
subsequently withdrawn. The separately refreshed canonical benchmark uses
261010; it is not a replacement label for the historical comparison.

\subsection{What Users Should Do}
\begin{enumerate}
\item \textbf{Check project inputs.} Low-burn forest should receive 0.06 m when
the revised shared template is selected and not overridden. Inspect prepared
\texttt{p*.man} files, not just a class label, displayed table or binary version.
\item \textbf{Regenerate deliberately.} Existing prepared files, project-local
managements and exported lookups are not automatically migrated. Regenerate
intended management artifacts and verify values before rerunning. Preserve
user-specified overrides.
\item \textbf{Keep passes compatible.} Watershed reruns require fresh hillslope
PASS files from the same WEPP build as the watershed executable.
Do not reuse cross-version passes.
\item \textbf{Interpret sediment in context.} Lower low-burn delivery may follow
this correction where effective roughness persists. Unchanged hillslope runoff
does not prove unchanged outlet sediment or require every small event to
deliver sediment.
\end{enumerate}

The extended lookup's four affected texture rows were updated; the base lookup
has no RRINIT column. Custom managements and generic initial-condition overrides
retain their precedence. Unrelated template/export discrepancies and regional
mappings were not normalized by this revision.

Representative observations or independent cases showing unacceptable
small-event sediment loss or other material regressions would justify review.
Rollback restores the shared template and all four lookup cells together;
it does not tune other parameters implicitly. This brief does not rerun projects
or authorize deployment.

\clearpage
\section{Decision and Evidence Record}
The accepted decision is ADR-0083. Implementation: WEPPpy commit
\par{\small\nolinkurl{bae734d71165e58d610d5282d41684113b18b4ff}}\par
Study binary SHA256:
\par{\small\nolinkurl{37d8deaf7a4c78e83104db4d896db3ee10b77a5f5d3f3abe5ad718390e2cc228}}\par

Implementation readback verified one template field and four lookup cells
changed: \texttt{ini.data.rrinit} in \nolinkurl{Low_Severity_Fire.man}
and the four corresponding rows in \nolinkurl{extended_land_soil_lookup.csv}.
All 16 generated low-forest managements parsed to 0.06 m and retained
0.04 m ridge height. Retained studies preserve setup failures, runtime
distinctions, controls and input/output receipts; they have not been silently
regenerated against today's defaults. The canonical climate and deterministic
generator are committed. No private watershed is a prerequisite for this
document's compilation or a mandatory publication gate.

For current benchmark expectations, see the
\href{https://wepp.cloud/weppcloud/usersum/src/tests/disturbed/analysis_results_current.md}{current canonical Disturbed rankings report}.
Check its version and input identities separately from this point-in-time brief.
